% Generated draft: requires mathematical and source review.
\documentclass[11pt]{article}
\usepackage[margin=25mm]{geometry}
\usepackage{fontspec}
\setmainfont{Latin Modern Roman}
\usepackage{amsmath,amssymb,mathtools,unicode-math}
\setmathfont{Latin Modern Math}
\usepackage{enumitem,needspace,xurl,hyperref}
\setlength{\parindent}{0pt}
\setlength{\parskip}{5pt}
\setlength{\emergencystretch}{3em}
\newcommand{\N}{\mathbb{N}}
\newcommand{\Z}{\mathbb{Z}}
\newcommand{\Q}{\mathbb{Q}}
\newcommand{\R}{\mathbb{R}}
\newcommand{\C}{\mathbb{C}}
\newcommand{\F}{\mathbb{F}}
\newcommand{\A}{\mathbb{A}}
\newcommand{\PP}{\mathbb{P}}
\newcommand{\E}{\mathbb{E}}
\newcommand{\eps}{\varepsilon}
\newcommand{\dd}{\,\mathrm{d}}
\begin{document}
\section*{UnsolvedMath ID: 700779}
\textbf{Source problem code:} catalogue entry
\subsection{Definitions and mathematical statement}
For each integer \(n\geq 4\), let \(F_n\) be the free group on \(n\) generators. For a group \(G\), write \(G'=[G,G]\) for its commutator subgroup, generated by the elements \([g,h]=g^{-1}h^{-1}gh\), and put \(G''=[G',G']\). A group is metabelian if its commutator subgroup is abelian. The free metabelian group of rank \(n\) is
\[
M_n=F_n/F_n''.
\]
Let \(\operatorname{Aut}(M_n)\) denote its group of all automorphisms, with composition as the group operation. A group is finitely presented if it has a presentation \(\langle X\mid\mathcal R\rangle\) with both \(X\) and \(\mathcal R\) finite.

The supplied source labels the following as a well-known problem: for every integer \(n\geq 4\), is \(\operatorname{Aut}(M_n)\) finitely presented? The finite presentation may depend on \(n\).
\subsection{Short English statement}
Is the automorphism group of the free metabelian group of rank \(n\) finitely presented for every integer \(n\geq 4\)? The supplied registry labels this problem open; that status has not been independently verified.
\subsection{Sources}
\begin{itemize}
\item Local UnsolvedMath registry, supplied record with permanent ID \(700779\), exact source problem code \texttt{catalogue entry}, title \emph{Kourovka Notebook Problem 11.87}. The supplied status is \texttt{open}.
\item \emph{The Kourovka Notebook (Unsolved Problems in Group Theory)}, Problem 11.87, attributed in the supplied record to V. A. Roman'kov, proposed in 1990, and described as a well-known problem. The background identifies the 11th issue (1990) and the 21st issue edition (2026), arXiv:1401.0300v48: \url{https://arxiv.org/abs/1401.0300v48}. The supplied collection description names E. I. Khukhro and V. D. Mazurov as editors. This bibliographic attribution is supplied source data and has not been independently checked.
\end{itemize}
\end{document}
